\begin{abstract}
I sistemi embedded distribuiti, come gli sciami robotici e le reti IoT,
richiedono un coordinamento autonomo e resiliente, senza un orchestratore
centrale, operando spesso con risorse computazionali e di memoria limitate.
\ac{AC} è un paradigma di macroprogrammazione adatto a questo scenario, ma molti
framework esistenti, come ScaFi, sono pensati per dispositivi general-purpose e
risultano incompatibili con i microcontrollori a basso costo.
\ac{YAAIR} nasce per superare questa limitazione, proponendo un'architettura a strati
che disaccoppia la logica del programma aggregato dal livello di rete, sfruttando
Rust per garantire prestazioni predicibili e un uso sicuro della memoria senza
compromettere sull'utilizzo delle risorse.

Il contributo principale di questa tesi consiste nella progettazione ed implementazione
di un livello di rete decentralizzato per \ac{YAAIR} basato sul protocollo Zenoh,
pensato per superare i vincoli topologici di soluzioni tradizionali come \ac{MQTT}
e per supportare dispositivi eterogenei, da embedded a sistemi general-purpose.

Il risultato finale è stato valutato tramite test sperimentali, volti a valutare
il corretto funzionamento della comunicazione tramite Zenoh tra dispositivi di
diverso tipo e l'integrazione della rete con \ac{YAAIR} per il suo utilizzo
nell'esecuzione di un programma aggregato.
\end{abstract}
